% !TEX root = ../main.tex

\chapter{Lampiran}

\section{Format Deklarasi Penggunaan AI}
\section{Surat Pernyataan Penggunaan AI}

\vspace{0.5cm}

\noindent\textbf{Dengan ini, kami:}
\begin{tabularx}{\textwidth}{@{}l X@{}}
\toprule
Nama / NIM & Andi Farhan Hidayat \& Naufarrel Zhafif Abhista / 13523128 \& 12523149 \\
Fakultas / Sekolah & Sekolah Teknik Elektro dan Informatika \\
Kode Mata Kuliah & IF3270 \\
Nama Mata Kuliah & Pembelajaran Mesin\\
Judul Tugas / Proposal & Implementasi \textit{Convolutional Neural Network} \& \textit{Recurrent Neural Network} \textit{From Scratch} \\
\bottomrule
\end{tabularx}

\vspace{0.5cm}

\noindent Menyatakan bahwa Kami menggunakan AI dalam pengerjaan maupun penyusunan luaran tugas pada mata kuliah yang tertulis di atas. 

\noindent Jika Ya, maka alat AI/kecerdasan buatan yang digunakan adalah:

\vspace{0.3cm}

\noindent\textbf{Detail Penggunaan}
\vspace{0.3cm}

\begin{tabularx}{\textwidth}{@{}p{6.5cm} p{2.5cm} X@{}}
\toprule
\textbf{Lingkup Pekerjaan} & \textbf{Digunakan?} & \textbf{Tingkat Penggunaan} \\
\midrule
Pemeriksaan Ejaan: menggunakan tools seperti Grammarly, Paperpal, DeepL, Quillbot, Grammarbot, LanguageTool, ProWritingAid, ChatGPT, Google Gemini, atau sejenisnya. & Tidak & 0 \\
\addlinespace
Pembuatan Teks: menggunakan tools seperti ChatGPT, Gemini, Paperpal, Copilot, GrammarlyGO, WordAI, WriteSonic, Jasper, Jenni AI, atau sejenisnya. & Ya & 3\\
\addlinespace
Bantuan Diskusi dalam Penyusunan Konten: menggunakan tools seperti ChatGPT, Gemini, Copilot, Perplexity, Jenni AI, Zoom-Companion, atau sejenisnya. & Ya & 2\\
\addlinespace
Pembuatan Gambar, Video, dan/atau Grafik: menggunakan tools seperti Craiyon, DALL\-E, Midjourney, Stable Diffusion, Microsoft Designer, Gemini, Canva AI, atau sejenisnya. & Tidak & 0 \\
\addlinespace
Lainnya (Jelaskan): & Ya/Tidak* & \\
\bottomrule
\end{tabularx}

\vspace{0.6cm}

\noindent Catatan: beri tanda \texttt{Ya} atau \texttt{Tidak} pada kolom "Digunakan?" dan jelaskan di "Lainnya" bila perlu. Tingkat penggunaan dibiarkan kosong untuk diisi kemudian.

\vspace{1cm}

\noindent\begin{tabularx}{\textwidth}{@{}X X@{}}
\centering (tempat), (tanggal) & 
\centering \\
\centering (tanda tangan) & 
\centering \\
\centering (Nama Lengkap) & \centering (NIM) \\
\end{tabularx}

\vspace{0.5cm}

\noindent Pemakaian tanda \texttt{*} menandakan opsi yang dipilih dapat diberi tanda centang sesuai kebutuhan.